\chapter{Modular Arithmetic}
\label{app:modular}

Modular arithmetic (sometimes ``clock arithmetic'') is ordinary
integer arithmetic --- addition and multiplication --- with one extra
step: after each operation, divide by the modulus and report only the
remainder.  For example, with modulus 5,
\[
  3 + 4 = 7 \equiv 2 \pmod 5,
\]
read ``three plus four is congruent to two modulo 5,'' since 5 goes
into 7 once with remainder 2.  Any integer reduces modulo 5 to one of
$0,1,2,3,4$: if positive, take the remainder on division by 5; if
negative, add 5 repeatedly until the result lies in that range.  Since
$5 \equiv 0 \pmod 5$, adding 5 never changes a result mod 5.  The
additive inverse of 1 is $-1$, and $-1+5 = 4$, so $-1 \equiv 4
\pmod5$.  In general
\[
  n \equiv 0 \pmod n, \qquad -1 \equiv n-1 \pmod n .
\]

Modulo-2 arithmetic has only two digits, $0$ and $1$, and reduces to
three rules (writing $=$ for congruence mod 2):
\[
  0 + 0 = 0, \qquad 0 + 1 = 1, \qquad 1 + 1 = 0 .
\]
These are exactly the rules of the ``exclusive or'' operation of
Boolean algebra, of central importance in the design of digital
circuits.  In Python, the operator is spelled \texttt{\%}
(\texttt{7 \% 5} evaluates to \texttt{2}), and it is applied
element-wise by NumPy: \texttt{np.mod(array, k)}.

\chapter{Pascal's Triangle}
\label{app:pascal}

The simplest construction of Pascal's triangle: start with a single 1;
each subsequent row is obtained by adding adjacent entries of the row
above (treating missing entries as 0):
\[
\begin{array}{ccccccccc}
 & & & & 1\\
 & & & 1 & & 1\\
 & & 1 & & 2 & & 1\\
 & 1 & & 3 & & 3 & & 1\\
1 & & 4 & & 6 & & 4 & & 1
\end{array}
\]
The entries are the \emph{binomial coefficients}, arising in the
expansion
\[
  (a+b)^n = \sum_{k=0}^{n} \binom{n}{k} a^k b^{n-k},
  \qquad \binom{n}{k} = \frac{n!}{k!\,(n-k)!},
\]
where $\binom{n}{k}$ sits at position $k$ of row $n$ (numbering rows
and positions from zero).

Reduced modulo 2, Pascal's triangle becomes an intricate pattern of
ones and zeros --- recognizable as none other than Sierpi\'nski's
gasket (Figure~\ref{fig:pascal2}).  This is exactly the pattern
produced by the additive CA with kernel $[1\;0\;1]$ modulo 2
(Section~\ref{sec:kernel101}), as the polynomial analysis of
Section~\ref{sec:poly1d} makes precise: the coefficients
$\binom{t}{i}$ appear directly in the closed form for that rule.
Computing $2^n$ rows of Pascal's triangle mod 2 yields a figure
resembling the approximation $S_n$ to the gasket
(Appendix~\ref{app:sierpinski}).

\begin{figure}[htbp]\centering
\includegraphics[width=0.7\textwidth]{appB_pascal_mod2}
\caption{Pascal's triangle modulo 2, 128 rows --- successive
approximations to Sierpi\'nski's gasket.}
\label{fig:pascal2}
\end{figure}

\chapter{Geometric Construction of Sierpi\'nski's Gasket}
\label{app:sierpinski}

In 1915 Wac\l{}aw Sierpi\'nski described the following geometric
construction of the fractal triangle now bearing his name.  Begin with
a filled equilateral triangle $S_0$.  To construct $S_1$, remove the
inverted middle triangle whose sides have half the length of the
original's.  To construct $S_2$, apply the same operation to each of
the three triangles composing $S_1$; to construct $S_3$, apply it to
each of the nine triangles composing $S_2$; and so on
(Figure~\ref{fig:sierp-construct}).

\begin{figure}[htbp]\centering
\includegraphics[width=\textwidth]{appC_sierpinski_construction}
\caption{The first four approximations $S_0, S_1, S_2, S_3$ to
Sierpi\'nski's gasket.}
\label{fig:sierp-construct}
\end{figure}

Continuing indefinitely yields a sequence $S_0, S_1, S_2, \dots$ of
better and better approximations; the gasket itself is the limit.  It
has zero area and Hausdorff dimension $\log 3 / \log 2 \approx 1.585$.
The CA with rule $[1\;0\;1]$ modulo 2, evolved for $2^{\,n+1}$
generations, produces a figure resembling $S_n$; the 256-generation run
of Figure~\ref{fig:sierp2} looks like $S_7$.  That three quite
different constructions --- a geometric removal process, Pascal's
triangle mod 2, and an additive CA --- converge on the same figure is a
small marvel worth pausing over.

\chapter{The Euler Totient Function}
\label{app:totient}

Euler's totient function $\varphi$ features prominently in number
theory.  For a positive integer $n$,
\[
  \varphi(n) = \text{the count of integers in } \{1,\dots,n\}
  \text{ that are relatively prime to } n,
\]
where two integers are \emph{relatively prime} if their only common
positive divisor is 1.  Thus 6 and 10 are not relatively prime (both
divisible by 2), while 9 and 10 are.  By convention $\varphi(1)=1$.
For a prime $p$ every smaller positive integer is coprime to it, so
\[
  \varphi(p) = p - 1 ,
\]
placing $\varphi$'s local maxima at the primes and its local minima at
integers with many factors (Figure~\ref{fig:totient-plot}).  More
generally, $\varphi$ is multiplicative and
$\varphi(p^a) = p^a - p^{a-1}$, giving the product formula
$\varphi(n) = n\prod_{p \mid n}(1 - 1/p)$.

\begin{figure}[htbp]\centering
\includegraphics[width=0.85\textwidth]{appD_totient}
\caption{Euler's totient function for $n < 200$.  Primes lie on the
upper line $\varphi(n)=n-1$.}
\label{fig:totient-plot}
\end{figure}

A short table:

\begin{center}
\begin{tabular}{r|l|r}
\toprule
$n$ & integers coprime to $n$ & $\varphi(n)$\\
\midrule
 2 & 1                         & 1\\
 3 & 1, 2                      & 2\\
 4 & 1, 3                      & 2\\
 5 & 1, 2, 3, 4                & 4\\
 6 & 1, 5                      & 2\\
 7 & 1, 2, 3, 4, 5, 6          & 6\\
 8 & 1, 3, 5, 7                & 4\\
 9 & 1, 2, 4, 5, 7, 8          & 6\\
10 & 1, 3, 7, 9                & 4\\
12 & 1, 5, 7, 11               & 4\\
\bottomrule
\end{tabular}
\end{center}

The module \texttt{cadyn.numtheory} computes $\varphi$ both singly
(\texttt{totient(n)}) and as a sieve over a range
(\texttt{totient\_table(limit)}); the latter is used by the
Euler-totient CA rule of Section~\ref{sec:totient}, with the convention
$\varphi(0)=0$ so that a dead (zero) cell contributes nothing to its
neighbors' sums.
